\documentclass[zihao=-4,a4paper]{ctexart}
\usepackage[margin=2.4cm]{geometry}
\usepackage{amsmath,amssymb,booktabs,graphicx,float,array,enumitem,xcolor}
\usepackage[most]{tcolorbox}
\usepackage[hidelinks]{hyperref}
\linespread{1.35}
\setlist{itemsep=2pt,topsep=4pt}
\ctexset{section={format=\Large\bfseries\color{blue!55!black}},subsection={format=\large\bfseries}}
\newtcolorbox{keybox}[1][本次要点]{colback=blue!3,colframe=blue!45!black,fonttitle=\bfseries,title=#1,boxrule=0.6pt,arc=1mm}
\newtcolorbox{notebox}{colback=gray!6,colframe=gray!60,boxrule=0.4pt,arc=1mm}
\newcommand{\fig}[3][0.95]{\begin{figure}[H]\centering\includegraphics[width=#1\linewidth]{#2}\caption{#3}\end{figure}}
\newcommand{\M}{M^{(K)}}

\title{\textbf{第 4 次组会汇报}\\[4pt]\Large 物理机理：组合风险从哪里来}
\author{对应工作：123、124 号实验；110—114、119 号的见证背景}
\date{}

\begin{document}
\maketitle
\thispagestyle{empty}

\begin{keybox}
\begin{enumerate}[leftmargin=*]
  \item 上一次的见证背景都能对上\textbf{已知的电力物理 / 市场机制}：比例恒等式、功率平衡与剩余负荷、阻塞与潮流、报价顺序。
  \item 受控机制算例（123 号）：\textbf{先写下机理预期，再做实验}。机组的私有报价加成由“本机组出力 + 本节点电价”反解；电价字段单独为 0，而 $M^{(2)}$ 把所有电价字段都识别为危险，见证背景都是本机组出力。其他方法全部失手。
  \item 字段风险画像（124 号）：同一字段在不同背景下的增益差异很大，$\M$ 是这个分布的上包络；\textbf{取平均会把背景放大稀释到约三分之一}。
  \item 危险背景往往“近等效”、共享一个核心伙伴字段，所以解释时应给出前几个背景，而不是只给最大的一个。
\end{enumerate}
\end{keybox}

\section{怎么读见证背景：增益矩阵}

增益矩阵的第 $i$ 行、第 $j$ 列表示：攻击者已经掌握字段 $j$ 时，再公开字段 $i$ 带来的推断增益 $V(\{i,j\})-V(\{j\})$。对角线（或第一列）是单字段能力。

\fig{figures/图1.png}{110 号：IEEE 30 节点的增益矩阵\protect\newline{\footnotesize 原图：\texttt{DNN\_Aggresvation110/figures/fig1\_gain\_matrix.png}}}

以目标 A（bus26 机组出力）为例：
\begin{itemize}[leftmargin=*]
  \item \textbf{阻塞标志 L24-26} 单独只有 0.01。但攻击者有风电预测时，它的增益是 0.23：知道风电多少、线路是否阻塞，就能判断这台机组是不是被线路“卡住”。
  \item \textbf{同一个风电预测}：攻击者已有阻塞信息时，公开它的增益是 0.69；已有系统备用时只有 0.21。\textbf{不说明攻击者手里有什么，字段风险就没有确定的值。}
  \item 负荷预测与实际负荷互为背景时增益约为 0：这是\textbf{冗余}，Shapley 会把它们对半分，$\M$ 不受影响。
\end{itemize}

\section{见证背景对应的四类物理机制}

\begin{table}[H]\centering
\caption{各场景见证背景的物理含义}
\footnotesize
\begin{tabular}{p{0.15\linewidth}p{0.36\linewidth}p{0.4\linewidth}}
\toprule
机制 & 例子（字段 + 见证背景 $\to$ 目标） & 物理含义 \\
\midrule
比例恒等式 & PJM 风电出力 + 风电占比 $\to$ 总发电（0.06 $\to$ 0.76，104 号） & 总发电 = 出力 $\div$ 占比 \\
\addlinespace
功率平衡 / 剩余负荷 & NEM 南澳风电预测 + 南澳可用容量 + 南澳光伏预测 $\to$ PPCCGT 出力（0.19 $\to$ 0.60，112 号） & 算出“南澳还需要燃气补足多少” \\
& NEM 昆士兰需求 + 昆士兰可用容量 + 新南威尔士—昆士兰联络线 $\to$ TUMUT3 出力（0.22 $\to$ 0.47） & 区域间功率平衡决定抽蓄要不要顶上 \\
& PJM 负荷预测 + 燃气出力 + 燃气占比 $\to$ 联络线净交换（0.07 $\to$ 0.52，113 号全候选口径） & 出力 $\div$ 占比 = 总发电，总发电 $-$ 负荷 = 净交换 \\
\addlinespace
净负荷 & CAISO 日前电能价 + 北部风电预测 + 南部光伏预测 $\to$ 实际负荷（0.14 $\to$ 0.51，114 号） & 电价加新能源预测反推净负荷位置 \\
\addlinespace
阻塞与潮流 & RTS 节点电价 + 光伏预测 $\to$ 线路 C35 潮流（约 0.5 $\to$ 0.7 以上，111 号） & 光伏决定潮流方向，电价反映阻塞状态，缺一不可 \\
& CAISO 实时 NP15 电价 + 实时 SP15 电价 $\to$ SP15 阻塞价（0.04 $\to$ 0.41） & 阻塞分量 $\approx$ 节点价差 \\
\addlinespace
报价顺序 & RTS 所在节点电价 $\to$ 机组出力（单字段即 0.90—0.995，111 号） & 电价决定机组在报价序列中的位置；10 个电价彼此冗余 \\
\bottomrule
\end{tabular}
\end{table}

\fig[0.8]{figures/图2.png}{111 号：RTS 线路 C35 的增益矩阵。电价彼此冗余（互为背景时增益接近 0），但配上光伏预测后各加 0.22—0.28\protect\newline{\footnotesize 原图：\texttt{DNN\_Aggresvation111/figures/fig2\_gain\_matrix.png}}}

\fig[0.8]{figures/图3.png}{112 号：NEM PPCCGT 的增益矩阵。“南澳可用发电容量”单独只有 0.02，却是几乎所有字段的最佳背景伙伴\protect\newline{\footnotesize 原图：\texttt{DNN\_Aggresvation112/figures/fig2\_gain\_matrix.png}}}

\begin{notebox}
PPCCGT 的例子说明，风险最大的\textbf{背景字段}本身可能单独毫无信息（南澳可用容量单独 0.02）。只看单字段能力，这个字段会被判为最安全，但它恰恰是打开其他字段的“钥匙”。
\end{notebox}

\section{受控机制算例：先写预期，再看能否找回（123 号）}

\subsection{设计}

前面的机理都是事后解释。为了验证 $\M$ 能找回\textbf{事先已知}的机制，在 RTS-GMLC 上构造：
\begin{itemize}[leftmargin=*]
  \item 网络、负荷、新能源、报价与 111 号完全相同，全年 8,784 小时 DC-OPF；
  \item 唯一改动：机组 313\_CC\_1（常为边际机组）的各段报价乘以 $(1+\kappa_t)$，$\kappa_t$ 是\textbf{私有的逐小时报价加成}，作为敏感目标。主算例 narrow 取 $\kappa\sim U(-5\%,5\%)$，加成小到基本不改变调度次序；
  \item 11 个字段，全部 2,047 个子集精确重训，所以任意 $K$ 都是精确值（$K=10$ 即完整 MCI）。
\end{itemize}

\textbf{事先写下的预期}：机组处于边际时 $\lambda_g=(1+\kappa_t)\,c_{\text{seg}}(P_g)$。出力 $P_g$ 定位报价段，电价 $\lambda_g$ 反解加成，所以 $\{P_g,\lambda_g\}$ 强协同；邻近节点电价在不阻塞时可以替代 $\lambda_g$；跨阻塞断面的远端电价在阻塞时失效。

\subsection{结果}

\fig{figures/图4.png}{123 号（narrow）：(a) 机组出力 $P_g$ 在各单背景下的增益；(b) 运行状态依赖；(c) 各方法打分（最后一列为 $M^{(2)}$）\protect\newline{\footnotesize 原图：\texttt{DNN\_Aggresvation123/figures/fig1\_mechanism\_narrow.png}}}

\begin{table}[H]\centering
\caption{123 号：各字段的风险（narrow）}
\small
\begin{tabular}{lcccl}
\toprule
字段 & $M^{(0)}$ & $M^{(2)}$ & 完整 MCI & $K=2$ 见证背景 \\
\midrule
机组出力 $P_g$ & 0.065 & \textbf{0.215} & 0.215 & 本节点电价 \\
本节点电价 & \textbf{0.000} & \textbf{0.150} & 0.150 & 机组出力 \\
邻近节点电价（306） & 0.000 & 0.141 & 0.141 & 机组出力 + C6 \\
远端电价（303，隔 C6） & 0.000 & 0.099 & 0.099 & 机组出力 \\
竞争机组出力 & 0.000 & 0.086 & 0.086 & 机组出力 \\
负荷、风光、阻塞 & 0.000 & 0.02—0.05 & $\le0.095$ & — \\
\bottomrule
\end{tabular}
\end{table}

\begin{itemize}[leftmargin=*]
  \item \textbf{与预期一致}：电价字段单独为 0，风险完全来自背景，见证背景都是本机组出力。$P_g$ 的增益排序为本节点电价 $>$ 邻近 $>$ 区域 2 $>$ 远端 $>$ 竞争机组 $\gg$ 负荷、风光、阻塞，与机理一致。
  \item \textbf{风险依赖运行状态}：C6 不阻塞时，远端电价 + $P_g$ 为 0.185，与本节点电价相当；C6 阻塞时两节点价格分离，远端电价降到 0.012，完全失效。
  \item \textbf{其他方法全部失手}（图 (c)）：相关系数、单字段、成对推断图给所有电价约 0；SAGE 把信用分摊到四个互相替代的电价上（本节点电价只剩 0.003）；LOCO 被替代字段互相掩盖；置换重要性只认本节点电价。
  \item wide 对照（$\kappa\sim U(-30\%,30\%)$）：加成大到改变调度次序时，出力单独就泄露（$M^{(0)}=0.524$），而电价仍是“单独为 0、背景下约 0.14”。\textbf{加成越小越隐蔽，单字段越看不见，但组合照样能推出来。}
\end{itemize}

\section{字段风险画像：背景放大有多普遍（124 号）}

定义\textbf{背景放大量} $\Gamma^{(K)}=M^{(K)}-M^{(0)}$，即最坏背景比单字段多出的风险。基于第 3 次的五组真值（10 个目标、264 个（字段，目标））。

\fig{figures/图5.png}{124 号：A 为单字段风险与背景放大量的关系（左上象限 = 单看不显眼、背景下被放大）；B 为同一字段在全部背景下的边际增益（空心 = 单字段，竖线 = 平均，实心 = 最大）\protect\newline{\footnotesize 原图：\texttt{DNN\_Aggresvation124/figures/fig1\_profile\_map.png}}}

\begin{table}[H]\centering
\caption{124 号：背景放大统计（264 个（字段，目标））}
\begin{tabular}{lc}
\toprule
指标 & 数值 \\
\midrule
$\Gamma^{(1)}$ / $\Gamma^{(2)}$ / $\Gamma^{(3)}$ 平均 & 0.082 / 0.110 / 0.119 \\
$\Gamma^{(2)}$ 单个最大 & 0.416 \\
第一个背景字段贡献的放大比例 $\Gamma^{(1)}/\Gamma^{(2)}$ & 75\% \\
单字段风险 $<0.3$ 而 $\Gamma^{(2)}>0.1$ 的条目 & \textbf{91 / 264} \\
平均背景增益 / 最大背景增益（中位数） & \textbf{32\%} \\
\bottomrule
\end{tabular}
\end{table}

\begin{itemize}[leftmargin=*]
  \item \textbf{$\M$ 是边际增益分布的上包络}：图 B 中增益分布很宽，最大值往往远离大多数背景。Shapley、SAGE 这类“平均语义”只能看到竖线的位置，大约是最大值的三分之一。
  \item \textbf{$K=1$ 有闭式解}：$\Gamma^{(1)}=\max\big(0,\max_j I_{ij}\big)$，其中 $I_{ij}=V(ij)-V(i)-V(j)$，在全部 264 个条目上误差为 0。也就是说，\textbf{一个背景带来的放大，就是字段与最强伙伴之间的正协同}。
  \item \textbf{危险背景近等效}：PPCCGT 的南澳风电预测，前 5 个背景的增益在 0.545—0.603 之间几乎并列，且都含“南澳可用发电容量”。所以解释输出应给前几个背景或它们共同的核心伙伴。
\end{itemize}

\section{什么时候第三个背景仍有用（119 号）}

\fig{figures/图6.png}{119 号：NEM 三台机组的 $K$ 递进（规模 $\le4$ 正式口径）\protect\newline{\footnotesize 原图：\texttt{DNN\_Aggresvation119/figures/fig2\_nem\_k3.png}}}

BW01、PPCCGT 的 $M^{(2)}/M^{(3)}$ 为 0.976 和 0.958，而抽蓄机组 TUMUT3 只有 \textbf{0.813}：它的见证背景含区域间联络线潮流。抽蓄要不要顶上，需要同时知道本区需求、可用容量和外区送入，这\textbf{天然是一个三方关系}。这也是论文中“$K=2$ 定级、$K=3$ 处置”的物理原因。

\section{小结与下一次}

\begin{keybox}[小结]
\begin{itemize}[leftmargin=*]
  \item 组合风险不是统计巧合：见证背景都对应比例恒等式、功率平衡、净负荷、阻塞潮流等已知机制。
  \item 在事先写下机理的受控算例中，$M^{(2)}$ 找回了全部预期结构，其他七种方法都没有找回。
  \item 背景放大普遍（35\% 的条目单看低风险、背景下被明显放大），取平均会稀释到约三分之一；最危险的“钥匙字段”可能单独毫无信息。
\end{itemize}
\end{keybox}

\textbf{下一次}：回到第 3 次留下的两个估计器问题。通用模型在真实数据上为什么会高估 $M$？消融实验发现\textbf{预学习几乎没用}，原因是什么，改成什么样才真正有用？

\end{document}
